\section{A uniform Morita reduction at rank ten}
\label{sec:rank10-morita}
Put $\delta=1+\sqrt2$. The five inputs
\[
 R10\_1\_114,\quad R10\_1\_115,\quad R10\_1\_152,
 \quad R10\_1\_153,\quad R10\_1\_179
\]
have dimension type $(1^4,(2\delta)^2,(\delta^2)^4)$.
We prove their noncategorifiability over $\C$. The finite computations in
this section reconstruct necessary module actions and fusion tensors;
they are not searches for numerical associators. The full source-bound
derivation and receipts are in
\href{https://github.com/sebastienpalcoux/grothendieck-ring-classification/tree/main/evidence/rank10_m1_exclusions_20261023}{the rank-ten continuation package}.

\subsection{The involution algebra and the forced dual}
In each input there is a unique nontrivial invertible $h$ fixing the two
simples $Y_0,Y_1$ of dimension $2\delta$. The remaining invertibles exchange
them, and $h$ acts freely on the other six simples. Directly in all five
tensors,
\begin{equation}\label{eq:rank10-odd}
 N_{Y_i,Y_j}^{Y_k}=1\qquad(i,j,k\in\{0,1\}).
\end{equation}
If $\mathcal C$ categorifies the input, the rank-one action of
$\langle h\rangle$ on $Y_0$ trivializes its pointed associator. Thus
$A=1\oplus h$ is a connected separable algebra. Its simple right modules
are four free modules from length-two $h$-orbits and two modules
$Y_i^+,Y_i^-$ over each fixed $Y_i$. Twisting the right action by the sign
automorphism of $A$ fixes the free modules and exchanges each split pair.

Let $W$ be the anti-invariant two-space with basis
$w_i=[Y_i^+]-[Y_i^-]$. Left multiplication induces a based-ring
representation $\rho$ on $W$. Left $h$ acts by
$\operatorname{diag}(\epsilon_0,\epsilon_1)$, with $\epsilon_i=\pm1$.
An invertible exchanging the $Y_i$ acts by an invertible signed swap.
Commutation with $h$ gives $\epsilon_0=\epsilon_1$. If both signs were
negative, $hY_i=Y_i$ would imply $\rho(Y_i)=0$. This is impossible:
each entry is $a-b$ with $a+b=1$ by \eqref{eq:rank10-odd}. Hence left
$h$ fixes every simple right module.

An $A$-bimodule is a right module with coherent left $C_2$ action.
Every right-module simple is fixed and has two one-dimensional equivariant
structures, since $H^2(C_2,\C^\times)=0$. The dual therefore has sixteen
simples of dimensions
\begin{equation}\label{eq:rank16-type}
 1^4,\quad\delta^8,\quad(\delta^2)^4.
\end{equation}
We use the endofunctor-composition convention
$\mathcal D=\operatorname{End}_{\mathcal C}(\mathcal C_A)$.
The corresponding literal bimodule convention takes the opposite tensor
product; all arguments below are invariant under this change.

Here the full eight-simple module action is forced, rather than just its
dimension vector. Normalize the signed swap on $W$. Its $+1$ eigenspace
has integer character values $y$ on both $Y_i$ and $z$ on the four
$\delta^2$-dimensional simples. The input tensor gives
\[
 y^2=2+2y+2z,\qquad yz=y+4z.
\]
Thus $(y-2)(y^2-4y-4)=0$, so $y=2,z=-1$. The odd entries of
$\rho(Y_i)$ are consequently all $+1$. Forgetting module structures
then determines the remaining entries; the fixed-sector matrices of
each $\delta^2$-dimensional simple are the negatives of the corresponding
identity or swap matrices. All module fusion and duality equations are
checked exactly by \nolinkurl{module_action.py}.

For clarity, the subsequent finite step has a proved bound. A simple
dual action matrix $B$ is nonnegative integral, commutes with the original
module action, and has prescribed unit column $Be_0=e_m$. Let $w$ be the
positive module dimension vector, with entries in $\{1,\delta,\delta^2\}$
and $w_0=1$. The sum of the original module matrices is symmetric and
irreducible. Its unique Perron line gives both
$Bw=\lambda w$ and $w^tB=\lambda w^t$, and the unit column gives
$\lambda=w_m$. Consequently
\[
 0\leq B_{ij}\leq\lambda w_i/w_j\leq\delta^4
 =17+12\sqrt2<34.
\]
The rational commutant has dimension twelve. Prescribing a unit column
leaves four parameters, each explicitly an entry of $B$. Necessary linear
inequalities narrow their integer ranges; exhaustive enumeration leaves
exactly two matrices for each of the eight columns. Tensoring with the
canonical sign bimodule $\tau$ exchanges these two matrices without fixed
points. Therefore both candidates occur as the two actual equivariant
structures; duplicate action matrices are not being assumed distinct.

Matrix products determine 240 of the 256 ordered dual products uniquely.
The remaining sixteen are precisely products of the four
$\delta^2$-dimensional simples. Each such simple is already a simple
product of two $\delta$-dimensional objects; associativity forces its
remaining products. The receipts verify all 4096 associativity equations
and all 256 matrix-action equations per dual tensor. This is a complete
necessary reconstruction with no arbitrary cutoff. The realized ring
$R10\_1\_21$ is included as an independent control.

\subsection{Relative factorization with a common invertible}
The adjoint subcategory $\mathcal B=\mathcal D_{\rm ad}$ has rank eight.
Its forced tensor has two rank-four fusion subcategories
$\mathcal P,\mathcal Q$ whose intersection has simples $1,f$.
The involution $\tau$ is outside $\mathcal B$, is distinct from $f$,
and conjugation by $\tau$ exchanges the two subcategories.
Every cross product of simples is simple; fixing either factor gives
four distinct products. Both rank-four factors have trivial universal
grading, and $f$ acts freely on their simple bases.

For inputs 114 and 115 their based rings are the ordinary
$\mathrm{PSU}(2)_6$ ring $P$:
\[
 f^2=1,\quad fx=y,\quad
 x^2=y^2=1+x+y,\quad xy=f+x+y,
 \qquad x^*=x,\ y^*=y.
\]
For inputs 152,153,179 they are the pseudo ring $Q$:
\[
 f^2=1,\quad fx=y,\quad
 x^2=y^2=f+x+y,\quad xy=1+x+y,
 \qquad x^*=y.
\]
All these identifications, the simple cross products, and the action of
$\tau$ are checked coefficient by coefficient in the forced tensors.

\begin{lemma}\label{lem:relative-factor}
In this situation multiplication extends to a dominant tensor functor
$\mathcal P\boxtimes\mathcal Q\longrightarrow\mathcal B$.
The unique central lifts of $f$ in the two factors have self-braidings
$q_{\mathcal P},q_{\mathcal Q}$ satisfying
$q_{\mathcal Q}=q_{\mathcal P}^{-1}=q_{\mathcal P}$.
\end{lemma}
\begin{proof}
Choose a noninvertible $Y\in\mathcal Q$. The complement
$\mathcal M=\mathcal P Y$ of $\mathcal P$ in $\mathcal B$ is a
$\mathcal P$-bimodule subcategory. The functor $X\mapsto XY$ is an
equivalence of left module categories. Transporting the right action
therefore identifies $\mathcal M$ with a twisted regular bimodule
$\mathcal P_\phi$. The tensor autoequivalence $\phi$ fixes simple
classes because $XY$ and $YX$ are the same simple class.

Multiplication followed by projection to $\mathcal M$ is a bimodule
functor $\mathcal P_{\phi^2}\to\mathcal P_\phi$. Its value at the
unit is $V=1\oplus f$, since the $\mathcal M$-part of $Y^2$ is
$Y\oplus fY$. Its right-module constraint gives coherent isomorphisms
$\phi^2(X)V\simeq V\phi(X)$, or, replacing $X$ by $\phi^{-1}(X)$,
$\phi(X)V\simeq VX$. On a simple $X$ the two constituent classes
$X,fX$ are distinct. These maps are therefore block diagonal for
$V=1\oplus f$. The unit block is a tensor isomorphism
$\phi\simeq\operatorname{Id}$. Thus $Y$ has an ordinary relative
half-braiding with $\mathcal P$.

Two relative half-braidings on a fixed simple $Z\in\mathcal Q$ differ
by scalars $\lambda_X$, because $XZ$ is simple. Tensor compatibility
makes them a character of the universal grading of $\mathcal P$,
which is trivial. The half-braiding is unique. Products of lifted
$\mathcal Q$-simples inherit half-braidings, and the latter preserve
each simple summand: tensoring by any fixed simple of $\mathcal P$
sends the four $\mathcal Q$-simples to distinct simples. Tensor powers
and summands of $Y$ generate $\mathcal Q$. These unique lifts are
therefore tensor closed and give the required functor into the relative
center, hence the dominant multiplication functor.

The relative crossing of the common $f$ with itself gives its unique
central lift in $\mathcal P$; reversing the crossings gives its lift
in $\mathcal Q$. Hence the two self-braidings are inverse.
Conjugation by $\tau$ identifies these central lifts, so they are also
equal. Their squares lie over the tensor unit and are trivial because
$\operatorname{Aut}_{\otimes}(\operatorname{Id})$ is trivial in each
factor. Thus $q_{\mathcal P}^2=1$.
\end{proof}

The case $q=+1$ is already impossible in either factor. It would give
the central boson algebra $1\oplus f$, whose fusion category of modules
has two simples and rule $W^2=1+2W$, excluded by the rank-two
classification \cite{Ostrik2}. The forced central lift is consequently
fermionic.

\subsection{The pseudo branch}
\begin{theorem}
The based rings $R10\_1\_152$, $R10\_1\_153$ and $R10\_1\_179$
have no complex fusion categorification.
\end{theorem}
\begin{proof}
Their rank-four factor has based ring $Q$ and, by
Lemma~\ref{lem:relative-factor}, a central fermion $f$. The standard
associated superfusion construction \cite[Definition~4.9]{Usher}
has morphism spaces
\[
 \operatorname{Hom}_{\widehat{\mathcal Q}}(U,V)_{\bar0}
 =\operatorname{Hom}_{\mathcal Q}(U,V),\qquad
 \operatorname{Hom}_{\widehat{\mathcal Q}}(U,V)_{\bar1}
 =\operatorname{Hom}_{\mathcal Q}(U,fV).
\]
The central self-braiding $-1$ gives exactly the super interchange sign.
Up to homogeneous isomorphism there are two simple representatives
$I=1,W=x$, both of m-type. The fusion spaces have dimensions
\begin{align*}
 \dim\operatorname{Hom}(W^2,I)&=(0\mid1),&
 \dim\operatorname{Hom}(W^2,W)&=(1\mid1),\\
 \dim\operatorname{Hom}(IW,W)&=(1\mid0),&
 \dim\operatorname{Hom}(WI,W)&=(1\mid0).
\end{align*}
Thus $W^2=\Pi I+(1+\Pi)W$. This is exactly the odd-evaluation family
$\widehat{\mathcal C}^{(1,1)}_m$ excluded by the finite superpentagon
calculation in \cite[Section~3]{CCX}. That calculation explicitly
includes the pair $(1,1)$ and admits only $(0,0)$ among the tested
odd-evaluation families. We use this finite result, not the subsequent
conjecture about all rank-two superfusion categories.

Neither the associated-superfusion construction nor that finite
polynomial-system exclusion assumes unitarity. In particular the cited
paper retains odd $(0,0)$ solutions even while discussing their failure
of a proposed physical locality condition; locality is not used as an
exclusion filter here. The final published Mathematica/Magma calculation
is a literature input and was not independently rerun in this project.
The precise rule and its applicability have been independently audited.
\end{proof}

\subsection{The ordinary branch}
\begin{theorem}
The based rings $R10\_1\_114$ and $R10\_1\_115$ have no complex
fusion categorification.
\end{theorem}
\begin{proof}
The earlier tensor-specific sphericalization and character receipts give
a pseudo-unitary Galois conjugate of any hypothetical categorification.
Pseudo-unitarity passes to its Morita dual: both global dimension and
Frobenius--Perron dimension are preserved
\cite[Theorem~2.15, Corollary~8.14]{ENO}; the positive spherical
structure then exists by \cite[Proposition~8.23]{ENO} and restricts to
the factors. Etingof--Ostrik's classification
\cite[Theorem~A.3]{EOSemisimple} identifies each pivotal ordinary
rank-four factor with $P=\operatorname{Rep}(\mathrm{SO}(3)_q)$,
$q$ a primitive sixteenth root. Conjugation by $\tau$ identifies the
two factors monoidally.

Lemma~\ref{lem:relative-factor} and central-algebra reconstruction
\cite[Corollary~5.6]{BrugNat} identify
$\mathcal B$ with $(P\boxtimes P)_a$, where
$a=1\oplus(f,f)$. The induced central structure is unique because
$U(P\boxtimes P)=1$. The common $f$ in each quantum factor is the
central fermion, so $(f,f)$ is a boson. The canonical inverse Morita
algebra is $1\oplus\tau$; its pointed associator is trivial. Hence
$\mathcal D$ is the crossed product of $\mathcal B$ by the coherent
order-two action $T=\operatorname{Ad}_{\tau}$, and the original
category is $\mathcal B^T$ (or its opposite under the opposite dual
convention).

We must exhaust the categorical actions, not just their based
permutations. Identify the second factor with the first using $T$.
The lift of $T$ to $P\boxtimes P$ is then ordinary factor interchange
$L$. Uniqueness of the cross-half-braidings gives a tensor comparison
$TF\simeq FL$, where $F:P\boxtimes P\to\mathcal B$ is induction.
Indeed applying $T$ to the crossing gives its inverse after exchanging
the factors, and these relative half-braidings are unique. Passing to
right adjoints gives a central algebra isomorphism $L(a)\to a$.
Conversely that isomorphism descends $L$ to $a$-modules: transport the
module action along it. Every module is a summand of a free one, so
the comparison recovers the action on every module and morphism.

The unital algebra automorphisms of $a$ are exactly the two signs on
$(f,f)$. Their squares are trivial. There is no extra action coherence,
since $U(\mathcal B)=1$ implies
$\operatorname{Aut}_{\otimes}(\operatorname{Id}_{\mathcal B})=1$.
Equivariantization commutes with this central condensation. The two
possible categories therefore have the form
\[
 ((P\boxtimes P)^{\rm swap})_{1\oplus(f,f)_{\epsilon}},
 \qquad\epsilon=\pm.
\]

For $\epsilon=+$, write $d_x^+$ for the positive equivariant lift of
$(x,x)$ and $p_{uv}$ for the off-diagonal orbit on $(u,v),(v,u)$.
In the permutation equivariantization,
\[
 (d_x^+)^2=d_1^++d_x^++d_y^++p_{1x}+p_{1y}+p_{xy}.
\]
Induction identifies $d_x^+,d_y^+$ and their common image $b$ is
simple, so $b^2$ contains $2b$. This contradicts multiplicity one.

For $\epsilon=-$ it suffices to determine the invertibles. Let
$s=F(d_1^-)$ and let $g,sg$ be the two summands of $F(p_{1f})$.
The pointed subcategory of $P\boxtimes P$ has trivial associator.
With $e_1=(1,f),e_2=(f,1)$, the equivariant structure exchanges the
two lines and the right $(f,f)_-$ action can be written
\[
 r(e_1)=ie_2,\qquad r(e_2)=-ie_1.
\]
The half-braiding is $-1$ on both lines, so the left action used in
the balanced product is $-r$. The quotient by the image of
$r\otimes1+1\otimes r$ is represented by its zero eigenspace, with
homogeneous lines
\[
 u=e_1e_1+e_2e_2,\qquad v=e_1e_2-e_2e_1.
\]
Their equivariance signs are $+,-$, respectively. The remaining right
action sends $u\mapsto iv$, $v\mapsto-iu$. Rescaling $v$ by $i$
identifies this module with the regular algebra. Hence $g^2=1$;
also $s^2=1$, and $s$ exchanges the two module structures. The four
invertibles form $C_2\times C_2$, whereas both target rings have
invertible group $C_4$.

Every primitive sixteenth root is Galois conjugate to the standard one.
The quantum category, central fermion, swap action, and algebra transport
under that conjugation. Fusion multiplicities and the invertible group
are unchanged. Thus the two contradictions apply to all forms of the
ordinary factors. Taking opposites does not alter either target based
ring, since both are commutative. This exhausts and excludes both inputs.
\end{proof}

\section{Structural constraints on the cyclic rank-ten ring}
\label{sec:cyclic-frontier}
The ring $R10\_1\_189$ has normal form
\begin{equation}\label{eq:cyclic-rank10}
 g^5=1,\quad gx=xg,\quad x^*=x,\quad
 x^2=1+(g+g^2+g^3+g^4)x,
\end{equation}
with $d=\FPdim(x)=2+\sqrt5$ and $D=50+20\sqrt5$.
The original basis map for
$(1,g,g^2,g^3,g^4,x,gx,g^2x,g^3x,g^4x)$ is
$(0,1,8,7,2,9,4,6,5,3)$. We record structural results valid for every categorification; the next section supplies an actual unitary model.

\begin{proposition}\label{prop:cyclic-associator}
Any complex categorification of \eqref{eq:cyclic-rank10} has pointed
subcategory tensor equivalent to the untwisted $\Vec_{C_5}$.
\end{proposition}
\begin{proof}
The established Galois reduction allows a positive spherical structure.
Vanishing of the pointed associator class is Galois invariant, so it
suffices to work there. Put $X_i=g^ix$, with subscripts modulo five.
The restriction Gram matrix has blocks
\[
 B_{GG}=10I,\qquad B_{GX}=5(J-I),\qquad B_{XX}=15I+15J.
\]
Writing a central restriction as $(a_i;b_i)$, the Galois automorphism
$\zeta_5\mapsto\zeta_5^2$ gives the exact transform
\[
 a'_i=a_{3i},\qquad b'_i=\sum_j a_j-a_{3i}-b_{3i}.
\]
These identities follow from the normalized center character formula;
the complete character matrix and the identity $HT=\sigma(H)$ are
verified exactly in the structural receipt. In particular $T^4=1$.
Every constituent of $I(g^j)$ has positive Galois sign and restriction
$g^j+\sum_{i\in S}X_i$, with $S\subset\{i:i\ne j\}$.

For completeness $I(g^j)$ is multiplicity free with ten constituents.
For $j=0$ this follows from commutativity of the based ring and induction.
For $j\ne0$, the fifth powers of their twists coincide
\cite[Proposition~3.8]{Schopieray}. The indicator identities
$\nu_m(g^j)=0$ for $1\leq m\leq4$ distribute the total dimension
equally among the five phases, each of weight $D/5=2+4d$.
If a constituent restricts with multiplicity $a$ to $g^j$ and has
dimension $a+bd$, it contributes $a(a+bd)$ to that weight. Comparing
constant coefficients gives $\sum a^2=2$ in each phase. There are
therefore two constituents, both of multiplicity one, in each phase.

Suppose the pointed associator class $k\in H^3(C_5,\C^\times)$ is
nonzero. The cyclic indicator formula \cite[Section~3.3]{MignardSchauenburg}
gives a primitive fifth root for $\nu_5(g)$. Equality in the trace
triangle inequality forces every constituent of $I(g)$ to have
primitive twenty-fifth-root twist. Choose a cyclotomic automorphism
$\sigma$ of exponent $2$ modulo $25$, with odd exponent on its
$2$-primary part. The normalized center has anomaly one, so modular
Galois symmetry gives
$\theta_{\widehat\sigma Z}=\sigma^2(\theta_Z)$
\cite[Theorem~2.9]{RankFinite}. Its fourth iterate fixes each restriction
but raises a primitive twenty-fifth-root twist to its sixth power,
an operation of order five. Thus each restriction in $I(g)$ occurs
a multiple of five times.

The ten restrictions are consequently five copies each of complementary
subsets $S,S^c$ of $\{0,2,3,4\}$. Their four Galois transforms supply
all forty nontrivial pointed-induced columns. The ten columns in $I(1)$
have $d$-coefficients $0,4,1,1,1,1,3,3,3,3$, from the exact formal
codegrees. Their singleton and complementary-triple counts coincide
coordinatewise by $FI(1)$, and Galois multiplication by two makes these
four counts equal, hence all one. Their noninvertible Gram contribution
has diagonal entries five and off-diagonal entries three away from
coordinate zero.

Subtracting these known columns from $15I+15J$ leaves diagonal ten.
For the eight complementary patterns, an off-diagonal entry of the
residual matrix is as follows; an entry exceeding ten contradicts a
nonnegative principal $2\times2$ minor:
\begin{center}\begin{tabular}{ccl}
\toprule Representative $S$ & Entry & Conclusion\\\midrule
$\varnothing$ & -- & five central lifts, treated below\\
$\{4\},\{3\},\{3,4\},\{2\},\{2,4\},\{2,3\}$ & $12$ & determinant $-44$\\
$\{2,3,4\}$ & $15$ & determinant $-125$\\\bottomrule
\end{tabular}\end{center}
The exact verifier checks all patterns, and also all 35 possible
singleton-count compositions before imposing their Galois equality.
The empty/full pattern would give five distinct central invertibles
forgetting to the same $g$. This is impossible: the kernel of forgetting
central lifts is $\operatorname{Aut}_{\otimes}(\operatorname{Id})$,
which is trivial here. Indeed scalars $u$ on $g$ and $v$ on $x$ satisfy
$v^2=1$ and $v^2=u^jv$ for $j=1,2,3,4$; hence $u=v=1$.
All nonzero pointed associator classes are excluded.
\end{proof}

\begin{proposition}
The regular-algebra Morita dual of any categorification of
\eqref{eq:cyclic-rank10} has the same based ring.
\end{proposition}
\begin{proof}
Proposition~\ref{prop:cyclic-associator} supplies $A=\bigoplus_{i=0}^4g^i$.
Its right-module category has two simples, and its dual has dimensions
$1^5,d^5$. This uses $H^2(C_5,\C^\times)=0$ when splitting the
left equivariant structures; no multiplicity-one assumption on the dual
is needed. The invertibles act freely and transitively on its five
noninvertibles. The commutation automorphism is either the identity or
inversion. Inversion would force five equal noninvertible coefficients
in a square, contradicting their total four. The ring is commutative.
Since five is odd, choose a self-dual noninvertible generator $z$.
Its square is
\[
 z^2=1+\{a+b(g+g^{-1})+c(g^2+g^{-2})\}z,
 \qquad a+2b+2c=4.
\]
Up to automorphism there are four types:
$(a,b,c)=(4,0,0),(2,1,0),(0,2,0),(0,1,1)$.
The first contains the excluded rank-two ring $z^2=1+4z$.
For the middle two, a formal codegree divided by five has respective
minimal polynomial
\[
 t^4-15t^3+60t^2-90t+45,\qquad
 t^4-20t^3+100t^2-160t+80.
\]
Both are Eisenstein at five and have degree four. But a Morita-equivalent
center has all simple dimensions in $\Q(\sqrt5)$, since forgetting to
the original category writes them in the original dimension semiring.
The central objects of dimensions $D/f$ furnished by induction therefore
exclude these quartic codegrees. The final type is exactly the original
ring. These polynomial identities and the four-type exhaustion are
checked with rational arithmetic in the Morita receipt.
\end{proof}

The irrational braided generalized-near-group classification
\cite[Theorem~3.6]{SchopierayBraided} excludes a braided categorification
with noninvertible dimension $2+\sqrt5$. The two propositions above
do not exclude a nonbraided categorification with untwisted pointed
subcategory. Equality of the dual's based ring is not a tensor
equivalence, and the rank-two module category is not a rank-two fusion
quotient. These necessary structural constraints supplement the actual unitary construction in Section~\ref{sec:r189-model}.
